Beyond these aggregate error profiles, performance variability across sessions was not primarily explained by event density but was more consistent with differences in signal quality and the sensitivity of sparse-blink sessions to individual detection errors. Performance varied substantially across the 104 sessions, from an $F_1$ of 0.9944 to 0.4167, with a median of 0.8730 and an IQR of [0.7659, 0.9360]. Subject-level means likewise ranged from 0.9928 to 0.5856, with a median of 0.8350. However, the lowest-scoring sessions were not those with anomalously high blink counts. The bottom five sessions had median ground-truth counts of 0.9x the dataset median on Raja and only 0.3x on Cao2018, despite dataset medians of 526 and 1314, respectively. Their errors were also not uniformly dominated by missed events, because only 2/5 Raja sessions and 1/5 Cao2018 sessions were FN-heavy. This pattern argues against event density as the primary failure mechanism. In sparse-blink sessions, a small number of detection errors imposes a larger penalty on the per-session score. Ambiguity in blink definition and heterogeneity in artifact morphology may further increase this sensitivity across recordings~\citep{nystrm2024what,arpaia2025systematic}.
